\documentclass[10pt,twocolumn,a4paper]{ctexart}
\usepackage[margin=1.7cm,columnsep=0.65cm]{geometry}
\usepackage{amsmath,amssymb,bm}
\usepackage{booktabs}
\usepackage{graphicx}
\usepackage{microtype}
\usepackage{natbib}
\usepackage{xcolor}
\usepackage{hyperref}
\usepackage{enumitem}
\usepackage{fancyhdr}

\graphicspath{{../../artifacts/figures/}}
\setlength{\parindent}{2em}
\setlist{nosep,leftmargin=*}
\hypersetup{colorlinks=true,linkcolor=blue!55!black,
  citecolor=green!45!black,urlcolor=blue!65!black}
\pagestyle{fancy}
\fancyhf{}
\fancyhead[L]{Data-Centric Offline RL Interview Guide}
\fancyhead[R]{ShaneLiu04}
\fancyfoot[C]{\thepage}

\title{\vspace{-1.2em}\textbf{数据决定上限：探索数据、CQL 与 IQL 的面试指南}}
\author{ShaneLiu04\\
\small\url{https://github.com/ShaneLiu04}}
\date{}

\begin{document}
\maketitle
\thispagestyle{fancy}

\begin{abstract}
Offline RL 的面试通常从算法开始：CQL 的正则为什么能压低外推，IQL 的 expectile 为什么
不查询分布外动作。工程上更常见的追问是数据：同一个算法换一份数据集，回报可以从
“永远到不了目标”变成接近最优。本文以 CS224R 2025 Homework 3 的 PointMass 实验为
线索，把三件事放在一起讲：无监督探索（RND 与随机策略）决定了数据覆盖；离散动作
CQL 的 logsumexp 正则在 $\alpha=0$ 时退化为离线 DQN 并发生过估计；IQL 的 expectile
在数据只覆盖部分路径时如何选择上分位。每一节都给出可检查的日志量与追问。
\end{abstract}

\textbf{关键词：}Offline RL；Exploration Data；RND；CQL；IQL；Coverage；Pessimism

\section{两阶段设定}
作业把一次运行拆成两段。前一段是无监督探索：智能体只最大化内在奖励，环境奖励不进入
探索 critic。随机探索在 \texttt{learning\_starts} 之前每步都均匀采样动作；RND 探索在
$\varepsilon=0.2$ 之外执行探索 critic 的 argmax。IQL 探索 $2\times10^4$ 步，命令不带
\texttt{--offline\_exploitation}，利用阶段仍以 $\varepsilon=0.2$ 的策略继续写入缓冲区，
实际上是 offline-to-online 的设定\citep{lee2022offline2online}。CQL 探索 $10^4$ 步后
缓冲区冻结，属于严格的 offline RL。

一个值得在面试里讲的细节：2025 starter 的 CQL agent 把探索策略的采集代码注释掉了，
每一步都调用 \texttt{env.get\_optimal\_action}，即按全源最短路距离贪心的专家控制器。
所以 CQL 的 \texttt{--use\_rnd} 只改变 critic 的训练起点，不改变数据。发现方式很朴素：
三组 CQL 的状态密度图逐像素相同。先画数据、再调算法，是 offline RL 工程的第一条纪律。

这种设定是 ExORL 的缩小版：先用与任务无关的探索算法采数据，再在同一份数据上用不同
奖励做离线学习\citep{yarats2022exorl,laskin2021urlb}。ExORL 的主要结论是
\emph{改数据比改算法更有效}：在探索性数据上，连普通的 TD3 都能接近 offline 专用算法。

\section{RND：用预测误差近似计数}
RND 固定一个随机初始化的目标网络 $f^*(s)$，训练预测网络 $\hat f_\phi$ 在缓冲区状态上
拟合它，奖励为 $\lVert\hat f_\phi(s')-f^*(s')\rVert$\citep{burda2019rnd}。见过越多次的
状态预测误差越小，所以它是伪计数的一种平滑近似\citep{bellemare2016count,tang2017exploration}。
与 ICM 相比，RND 的预测目标是确定性的，不受环境随机性（noisy TV）影响
\citep{pathak2017icm}。实现细节：
\begin{itemize}
\item 奖励要做归一化（作业用滑动标准差），否则不同阶段的量级差几个数量级；
\item 目标网络与预测网络使用不同的初始化，避免一开始误差就接近 $0$；
\item RND 奖励是非平稳的：预测网络在学，奖励在衰减，探索 critic 必须持续更新。
\end{itemize}
最大熵探索给出另一种理论视角：理想的无监督探索最大化状态访问分布的熵
\citep{hazan2019maxent}。均匀随机策略在迷宫里的状态分布高度集中在起点附近，
因为扩散的半径只随 $\sqrt{t}$ 增长；RND 主动走向低访问区域，覆盖面更大。

\section{数据覆盖与可学性}
理论上，悲观的 offline RL 只需要数据覆盖最优策略的状态动作分布，而不需要覆盖全部
\citep{jin2021pessimism,rashidinejad2021bridging}。这给出一个可操作的判断：如果数据里
没有任何一条轨迹到达目标附近，任何离线算法都无法学出到达目标的策略；覆盖足够时，
算法的差别才显现。Schweighofer 等人用“轨迹质量”和“状态覆盖”两个指标刻画
D4RL 式数据集，并发现覆盖度高的数据上 offline RL 优于行为克隆，质量高、覆盖低的数据
上行为克隆更稳\citep{schweighofer2022dataset,kumar2022should}。

在 PointMass Medium 上，可以直接用作业输出的 \texttt{curr\_state\_density.png}
检查覆盖。IQL 的随机探索密度集中在起点附近，RND 扩展到迷宫远端，IQL 的回报相应
高出 $30$ 以上。我们让 CQL 真正使用探索数据后得到一组更有意思的对照
（Medium，各 $6$ 个种子）：专家数据上 $\alpha=0$ 全部失败（$-150$），$\alpha=0.1$ 约
$-68$，与专家控制器自身的 $-67$ 持平；均匀随机数据上 $\alpha=0$ 的离线 DQN 也能学到
约 $-52$，$\alpha=0.1$ 达到 $-34$，明显超过专家；RND 数据上 $\alpha=1$ 全部失败，因为
强保守把策略拉向偏向“新奇”的探索行为。结论：保守系数要与数据的行为策略一起选。

\section{离散动作 CQL}
离散动作下，CQL 的正则可以精确计算\citep{kumar2020cql}：
\begin{equation}
\mathcal{L}=\mathcal{L}_{\mathrm{TD}}+\alpha\,\mathbb{E}_{s\sim\mathcal{D}}\Bigl[
\log\textstyle\sum_a e^{Q(s,a)}-\mathbb{E}_{a\sim\mathcal{D}}Q(s,a)\Bigr].
\end{equation}
第一项是软最大值，压低所有动作；第二项抬高数据动作。最优点上，分布外动作的 Q 值
被压到数据动作以下，策略改进就不会选到它们。实现要点：
\begin{itemize}
\item logsumexp 用 \texttt{torch.logsumexp}，不要手写 \texttt{log(sum(exp))}；
\item 第二项取的是 $Q(s,a_{\text{data}})$，不是 $\max_a Q$；
\item 日志里同时记录 $\log\sum e^{Q}$ 与 $Q(s,a_{\text{data}})$，两者之差就是
“分布外间隙”，是判断正则是否生效的最直接指标。
\end{itemize}

\paragraph{$\alpha=0$ 意味着什么。}正则消失，剩下的是在固定缓冲区上训练的 double DQN。
数据只覆盖部分动作时，$\max_a Q(s',a)$ 会挑中没有被数据约束的动作，误差经备份传回
\citep{fujimoto2019bcq,kumar2019bear}。作业的专家数据上，$\alpha=0$ 的六个种子回报全部是
$-150$，即 episode 长度上限。数据动作的 Q 值约 $-1.0$，而每步奖励 $-1$、到达需要几十步，
显然是过估计；分布外间隙为 $1.91$，$\alpha=0.1$ 时降到 $0.48$，$\alpha=3$ 时为 $0.12$。

\paragraph{$\alpha$ 的取值。}$\alpha$ 过大时 Q 被压向行为策略的价值，CQL 退化为
加权的行为克隆，回报趋于数据质量；过小时外推重新出现。原论文还给出
Lagrangian 版本，让 $\alpha$ 自适应地把间隙维持在阈值 $\tau$\citep{kumar2020cql}。
后续工作指出 CQL 在数据覆盖好的区域过于保守，MCQ 只对分布外动作施加温和的
伪目标\citep{lyu2022mcq}，EDAC 用 Q 集成的不确定性代替显式正则\citep{an2021edac}，
Cal-QL 校准保守价值以便在线微调\citep{nakamoto2023calql}。

\section{IQL 与 expectile}
IQL 用 expectile 回归拟合 $V(s)$\citep{kostrikov2022iql,newey1987expectile}：
\begin{equation}
L_V=\mathbb{E}_{(s,a)\sim\mathcal{D}}\bigl[|\zeta-\mathbb{1}(u<0)|\,u^2\bigr],
\quad u=Q(s,a)-V(s),
\end{equation}
Q 的目标为 $r+\gamma(1-d)V(s')$，策略用 AWR 抽取\citep{peng2019awr,nair2020awac}：
$\exp(A/\lambda)$ 加权的对数似然。$\zeta=0.5$ 时 $V$ 是行为策略的期望价值，等价于
SARSA；$\zeta\to1$ 时 $V$ 逼近数据支撑内的最大 Q，接近 in-sample 的 Q-learning。
所有更新只用数据中出现过的动作，因此不会外推到分布外动作。SQL/EQL 把这一思路
推广为隐式价值正则\citep{xu2023sql}，XQL 用 Gumbel 回归达到类似效果
\citep{garg2023xql}。

\paragraph{为什么 $\zeta$ 不是越大越好。}$\zeta$ 越大，$V$ 越依赖少数高价值样本，
方差越大；在随机环境中还会对“运气好”的转移过度乐观。作业中 Easy 环境的
$\zeta$ 差别不大，因为数据已经覆盖到目标；Medium 环境中 $\zeta$ 的影响更明显。

\section{工程清单}
\begin{itemize}
\item 先看数据：状态密度图、到达目标的轨迹数、平均回报。覆盖不够时不要调算法。
\item CQL：监控分布外间隙与数据 Q 值。间隙趋于 $0$ 且回报差，说明 $\alpha$ 太小；
数据 Q 值远低于真实回报，说明 $\alpha$ 太大。
\item IQL：$V$ 损失与 Q 损失应同时下降；优势的尺度决定 AWR 权重，必要时截断
$\exp(A/\lambda)$。
\item 多种子：PointMass 的回报是“到达/未到达”的混合，三个种子的方差很大，要报告
逐种子结果\citep{agarwal2021precipice}。
\item 区分 offline 与 offline-to-online：缓冲区是否在利用阶段继续增长。
\end{itemize}

\section{追问清单}
\begin{enumerate}
\item 为什么同样的 CQL 在 RND 数据上更好？如果给随机探索 $10$ 倍步数呢？
\item 离散动作 CQL 的正则与 softmax 交叉熵有什么关系？（$\log\sum e^{Q}-Q(s,a)$
正是以 Q 为 logits 的交叉熵。）
\item $\alpha=0$ 的 CQL 与行为克隆谁更好？在什么数据上？
\item IQL 的 $\zeta$ 与 CQL 的 $\alpha$ 分别控制哪一类误差？
\item 如果允许在线微调，CQL 的保守价值会带来什么问题？\citep{nakamoto2023calql}
\item 为什么 RND 对 noisy TV 鲁棒？在什么环境里它仍然会失效？
\end{enumerate}

\bibliographystyle{plainnat}
\bibliography{../references}

\end{document}
